\documentclass[12pt,a4paper]{exam}

\usepackage[utf8]{inputenc}
\usepackage[T1]{fontenc}
\usepackage[french]{babel}
\usepackage{amsmath,amssymb}
\usepackage{pgfplots}
\pgfplotsset{compat=1.18}
\usepackage{tikz}
\usetikzlibrary{arrows.meta}
\usepackage{geometry}
\geometry{margin=2cm}

\setlength{\parindent}{0pt}

\begin{document}

\begin{center}
\Large \textbf{- Interrogation -\\Notions sur les fonctions (2nde)}
\vspace{0.5cm}
\\\large \underline{\textbf{RÉPONDRE ENTIÈREMENT SUR UNE COPIE SÉPARÉE}}
\end{center}
\vspace{0.4cm}
\noindent\hrulefill
\vspace{0.5cm}

% --------------------------- %
% ------- EXERCICE 1 -------- %
% --------------------------- %

\textbf{Exercice 1 – Cours et parité (6 points)}
\vspace{5pt}

\begin{center}
{\renewcommand{\arraystretch}{2.5}
\begin{tabular}{|c|p{4cm}|p{4cm}|}
    \hline
    \text{Fonction} & \multicolumn{1}{c|}{\text{Expression}} & \text{Ensemble de définition} \\
    \hline
    \text{Carrée} & \dotfill & \dotfill \\
    \hline
    \text{Cube} & \dotfill & \dotfill \\
    \hline
    \text{Racine carrée} & \dotfill & \dotfill \\
    \hline
    \text{Inverse} & \dotfill & \dotfill \\
    \hline
\end{tabular}
}
\end{center}
\vspace{5pt}

\begin{questions}
    \question[2] Compléter le tableau.
    \question[1] \underline{Question de cours :} Donner la définition algébrique et graphique d'une fonction
    paire.
    \question[2] On considère la fonction suivante :
    \[
    f(x) = x^3 -2x + 4
    \]
    \begin{parts}
        \part Calculer $f(-x)$.
        \part La fonction est-elle paire ? Impaire ?
    \end{parts}
    \question[1] Calculer f(2). Sans le calculer, que vaut alors f(-2) ?
\end{questions}

\noindent\hrulefill
\vspace{0.5cm}

% --------------------------- %
% ------- EXERCICE 2 -------- %
% --------------------------- %

\textbf{Exercice 2 – Résolutions (8 points)}
\vspace{5pt}

Résoudre algébriquement ou graphiquement les équations et inéquations suivantes. Écrire 
les résultats sous forme d'\textbf{ensemble de solutions} ($S = ...$).

\vspace{5pt}

\begin{questions}
    \question Résoudre l'équation : $x^2 = 3$
    \question Résoudre l'équation : $x^3 = -27$
    \question Résoudre l'équation : $\sqrt{x} = 2$
    \question Résoudre l'inéquation : $x^2 \leq 9$
    \question Résoudre l'inéquation : $\sqrt{x} < 3$
    \question Résoudre l'inéquation : $\dfrac{1}{x} \geq 2$
\end{questions}

\noindent\hrulefill
\vspace{0.5cm}

% --------------------------- %
% ------- EXERCICE 3 -------- %
% --------------------------- %

\textbf{Exercice 3 – Étude de courbes (6 points)}
\vspace{5pt}
\\On donne ci-dessous les courbes suivantes :
\[
\begin{array}{ccc}
    y = x \ \ ;
    & y = x^2 \ \ ;
    & y = x^3
\end{array}
\]

\begin{center}
\begin{tikzpicture}[scale=2]

% Grille
\draw[step=1, gray!40, thin] (-1,-1) grid (3,3);

% Axes
\draw[->, thick] (-1,0) -- (3,0) node[right] {$x$};
\draw[->, thick] (0,-1) -- (0,3) node[above] {$y$};

% Graduations en x
\foreach \x in {-1,0,1,2,3}
    \draw (\x,0.1) -- (\x,-0.1) node[below] {\small \x};

% Graduations en y
\foreach \y in {-1,0,1,2,3}
    \draw (0.1,\y) -- (-0.1,\y) node[left] {\small \y};

% Courbe y = x
\draw[domain=-1:3, samples=200, thick, blue]
    plot(\x,{\x});
\node[blue] at (3.5,3) {$y=x$};

% Courbe y = x^2
\draw[domain=-1:1.77, samples=200, thick, red]
    plot(\x,{\x*\x});
\node[red] at (2.2,3.2) {$y=x^2$};

% Courbe y = x^3
\draw[domain=-1:1.5, samples=200, thick, olive!70!black]
    plot(\x,{\x*\x*\x});
\node[olive!70!black] at (1,2.5) {$y=x^3$};

\end{tikzpicture}
\end{center}

\vspace{5pt}
On se place sur l'intervalle $\left[0;3\right]$
\vspace{5pt}
\begin{questions}
    \question[1] Comparer $x$ et $x^2$ sur $\left[0;1\right]$.
    \question[1] Comparer $x^2$ et $x^3$ sur $\left[1;3\right]$.
    \question[2] Résoudre \textbf{graphiquement} l'inéquation : $x > x^3$.
    \question[2] Résoudre \textbf{graphiquement} l'équation : $x = x^3$.
\end{questions}

\noindent\hrulefill
\vspace{0.5cm}

\begin{center}
\textbf{Fin de l'interrogation — relire avant de rendre.}    
\end{center}

\end{document}
